% Source: sources/returns/S6_04_MAPS_MANIFOLDS/sections/fourfolds_D2_c3.tex:1–234
\paragraph{Theorem D.2.39 (Proved: the full nine-valued invariant in the $c=3$ sector).}
Use the primitive Mumford filling at every multiplicative site in one
of the following eight minimal patterns, retain the indicated additive
choice and the ordinary clutch $e\gamma$, and make no additional
regular logarithmic transformations:
\[
\begin{array}{c|r|r|l|l}
\text{OS/profile}&a&b&\text{zero positions}&\text{additive choice}\\\hline
56&5&3&2&\text{none}\\
56&10&3&3,4&\text{none}\\
56/II&5&3&0&II\text{ original or log}\\
56/III&1&6&0&III\text{ log}\\
56/III&5&6&3&III\text{ log}\\
56/III&10&3&3&III\text{ original}.
\end{array}                                                   \tag{D.T31}
\]
Zero positions are numbered from zero. For every finite multiplicity
and affine twist producing a simply connected manifold, the exact
oriented diffeomorphism type is
\[
 M_{3,0}=\bigl((S^2\times S^4)\setminus
          \operatorname{int}(S^2\times D^4)\bigr)
               \mathop{\cup}_{S^2\times S^3}(D^3\times S^3),     \tag{D.T32}
\]
where the embedded sphere represents three times the generator of
$H_2(S^2\times S^4;\mathbb Z)$ and has its oriented normal framing.
For example it is the graph of the degree-three map
$z\mapsto z^3$ on $\mathbb {CP}^1=S^2$ and an equatorial embedding
$S^2\hookrightarrow S^4$.
The complete Zhubr data are
\[
 H_2=H_3=\mathbb Z/3,\qquad b_3=0,\qquad
 \mathfrak p=0,\qquad
 \mu=0\in H_6(K(\mathbb Z/3,2);\mathbb Z)=\mathbb Z/9.          \tag{D.T33}
\]
In particular the assertion about $\mu$ determines its order-three
subgroup as well as its ordinary mod-three cubic reduction.

\begin{proof}
\emph{1. The invariant which must be computed.}
We first identify its complete target from the primary classification
and the actual Pontryagin operation. Zhubr's epimorphism in Section~3.11
and his free calculation in Section~3.12 show that
$\Omega^{\mathrm{Spin}}_6(K(\mathbb Z/3,2))$ is generated by the
images of $\epsilon,\zeta$, where
\[
 (P,h)(\epsilon)=(6,0),\qquad
 \zeta=[\mathbb {CP}^3,\text{hyperplane}].
\]
The image of $\epsilon$ has both invariants zero: its first is six
modulo three, and its second is the image of zero. The injectivity
in Zhubr's Theorem~5 therefore kills that image.
Lemma~3.14 gives $N(\mathbb Z/3)=0$, so the fundamental-class map
identifies this cyclic bordism group with
$H_6(K(\mathbb Z/3,2);\mathbb Z)$; the mod-two obstruction target
is zero because all positive homology is three-primary.
These facts use \cite[Sections 3.11--3.14 and Theorem 5]{manZhubr75}.

The cyclic generator has order dividing nine. On
$\mathbb {CP}^1\times\mathbb {CP}^2$ take the integral class
$3x+y$, with $x,y$ the two hyperplane classes. Its mod-three
reduction factors through $\mathbb {CP}^2$, and hence sends the
six-dimensional fundamental class to zero. Before reduction its
cube evaluates to nine. Naturality therefore gives nine times the
universal generator equal to zero.
For the opposite bound, Thomas's Pontryagin third power is a natural
operation
\[
 \mathcal P_3:H^2(-;\mathbb Z/3)\longrightarrow
                              H^6(-;\mathbb Z/9).
\]
Its value on the reduced hyperplane class of $\mathbb {CP}^3$
evaluates to a generator of $\mathbb Z/9$
\cite[equation (2.8), p.267]{ffThomas56}. Thus the cyclic group has
order exactly nine. Evaluation of $\mathcal P_3$ detects its entire
generator, and hence the whole fundamental-class invariant. This is
stronger than evaluating the ordinary cube modulo three.

\emph{2. The common-piece comparison and its circle geometry.}
D.2.26 gives $H_2(X)=H_3(X)=\mathbb Z/3$ and $b_3=0$ for every
simply connected member in (D.T31). In particular $H^2(X;\mathbb Z)=0$,
so $c_1(X)=0$ and $X$ is spin. Let
$u\in H^2(X;\mathbb F_3)$ be the generator normalized below.
All the geometric hypotheses of the doubled-line calculation in
D.2.36 hold here. We record the exact inputs to ensure that no
prime-greater-than-three conclusion is being imported.

Remove the additive caps and one ordinary clutch disc to form $W$.
In all seven marked initial rows the actual multiplicative angular
columns generate $\ker\gamma^*$ with three unit Smith factors.
Consequently $\pi_1(W)=\mathbb Z_\gamma\times F_r$ is torsion-free.
The integral degree-one bridge calculation and surjective degree-two
cap boundary maps give $H_2(V)\twoheadrightarrow H_2(W)$, as in
(D.C26). The pure fibre functional is
$\Omega=a\,u_1^*u_2^*+b\,\gamma^*\delta^*$, with $3\mid b$
and $3\nmid a$. Its finite normal coefficients are divisible by
$b/q_i=3$, its original additive coefficients are zero, and the
ordinary clutch coefficient is $be$. Thus its mod-three reduction
is the actual global degree-two generator, with zero marked swept
values. The doubled factors (D.OP4) and the exact integral supports
(D.OP5), (D.OP5a) produce the corresponding line on $W$.

The scalar circle has precisely two strata whose stabilizer order is
divisible by three: the rational curve $C$ fixed by the whole circle,
and one elliptic double curve $E$ with stabilizer exactly $C_3$.
The latter is at the unique $I_3$ position. On its complement all
stabilizers are finite of order prime to three. This includes the
finite logarithmic caps: primitivity of their actual gamma numerators
makes the scalar circle free there, for every allowed order.
The original additive caps have free scalar action as well.
The seven exact Tate markings in the certificate verify the remaining
multiplicative assertions, including the primitive short period
$(3,aj)$ at $E$ and the two normal characters
$w,-w$, where $w=(aj)^{-1}=\pm1$ in $\mathbb F_3$.

The doubled auxiliary line has zero $C_3$ character at both $C$ and
$E$. Its degree on $C$ is a unit modulo three. Normalize $u$ by
$\langle u,[C]\rangle=1$. The torsion-free group of $W$ gives the
same simultaneous fixed-point comparison as in D.2.36: once an
equivariant lift has zero point character on $C$, it also has zero
point character on $E$. Indeed the two fixed-point copies of $BC_3$
represent the same order-three subgroup in
$\pi_1(W_{C_3})=\pi_1(W)\times C_3$. The group-cohomology difference
from half the auxiliary-line class has a common degree-two base
coefficient; the possible mixed degree-one term does not change
that coefficient. This comparison uses the complete normal columns,
including arbitrary finite twists.

\emph{3. The full-circle lift and the elliptic Borel space.}
The generator $u$ admits a lift
\[
                 \widehat u\in H^2_{S^1}(X;\mathbb F_3).
\]
In the Borel spectral sequence the only possible differentials out
of ordinary degree two have targets involving $H^1(X;\mathbb F_3)$
or $H^3(BS^1;\mathbb F_3)$; both are zero. Adding a multiple of the
base generator normalizes its restriction on the fixed curve to
\[
                      \widehat u|_{C_{S^1}}=[C]^* .           \tag{D.T34}
\]
Restriction $H^2(BS^1;\mathbb F_3)\to H^2(BC_3;\mathbb F_3)$
is nonzero. The point-character comparison in part~2 consequently
makes the coefficient of the degree-two $BC_3$ class on $E$ zero
as well.

The effective action $S^1/C_3$ on the elliptic curve $E$ is free.
Its quotient is a circle. Every principal circle bundle over a
circle is trivial, so equivariantly $E$ is the product of a free
$S^1/C_3$ orbit and a circle. Therefore
\[
 E_{S^1}\simeq BC_3\times S^1,
 \qquad
 \widehat u|_{E_{S^1}}=\eta\times\zeta,                    \tag{D.T35}
\]
where $\eta\in H^1(BC_3;\mathbb F_3)$ and
$\zeta\in H^1(S^1;\mathbb F_3)$, allowing a scalar in either
factor. Formula (D.T35) retains the possible nonzero mixed class;
it does not assert that $u$ descends to the coarse quotient $X/C_3$.
In fact the mixed term is generally essential.

For odd primes Thomas proves exact external-product multiplicativity
and vanishing on odd-dimensional classes, and independently proves
vanishing on suspensions
\cite[equations (2.3), (2.5), (2.7), p.267]{ffThomas56}.
These are statements with the full target coefficients, so
\[
 \mathcal P_3(\widehat u)|_{E_{S^1}}
   =\mathcal P_3(\eta)\times\mathcal P_3(\zeta)=0
                       \quad\text{in }H^6(E_{S^1};\mathbb Z/9).
                                                               \tag{D.T36}
\]
There is also no contribution on $C_{S^1}$: (D.T34) is pulled back
from the two-dimensional curve, whose degree-six cohomology is zero.
The familiar nonvanishing phenomena for the Pontryagin square of
suspended classes are irrelevant here; the cited formulas explicitly
assume an odd prime.

\emph{4. Why the full order-nine class vanishes globally.}
Put $R=\mathbb Z/9$. Choose disjoint invariant tubular neighborhoods
$N$ of $C\sqcup E$ and let $Y=X\setminus\operatorname{int}N$.
The circle acts locally freely on $Y$, with stabilizer orders prime
to three. The quotient $Y/S^1$ is a compact five-dimensional orbifold
with boundary. The Borel projection to this quotient, for both $Y$
and its boundary, has local cohomology fibres $BH$ with $H$ finite
of order invertible in $R$. Finite averaging gives
$H^q(BH;R)=0$ for $q>0$. The Leray comparison and the long exact
sequences of the pairs therefore give
\[
 H^k_{S^1}(Y,\partial Y;R)
             \cong H^k(Y/S^1,\partial(Y/S^1);R)=0
                            \qquad(k>5).                     \tag{D.T37}
\]
The dimension bound can be read from an orbifold triangulation,
or from a finite good orbifold atlas and its relative cellular
cohomology. No division by three occurs in this comparison.

Equivariant excision identifies the left side of (D.T37) with
$H^k_{S^1}(X,N;R)$. At $k=6$ it is zero. Thus restriction
\[
 H^6_{S^1}(X;R)\longrightarrow H^6_{S^1}(C\sqcup E;R)
                         \text{ is injective}.               \tag{D.T38}
\]
The vanishing of both restrictions proved in part~3 now implies
$\mathcal P_3(\widehat u)=0$ in the full integral-coefficient group
with coefficients $R$. Restricting to the ordinary fibre $X$ gives
\[
                        \langle\mathcal P_3(u),[X]\rangle=0
                                     \quad\text{in }\mathbb Z/9.
                                                               \tag{D.T39}
\]
Part~1 proves that this is the entire invariant $\mu$. Since
$N(\mathbb Z/3)=0$, the spin-bordism class itself is zero, and its
refined Pontryagin class is zero as well. This also agrees with
ordinary localization, which yields $p_1=0$ because $w^2=1$;
that latter congruence alone would not have proved (D.T39).

\emph{5. The stated concrete smooth model.}
The graph embedding in (D.T32) represents three times the generator.
Its oriented real rank-four normal bundle is trivial: the ambient
manifold is spin and $w_2(TS^2)=0$, and $w_2$ classifies such bundles
on $S^2$. The framing is unique up to homotopy because
$\pi_2SO(4)=0$. The framed surgery preserves simple connectivity
and spin. Its integral sequence and duality give
$H_2=H_3=\mathbb Z/3$, with $b_3=0$.
The mod-three class from the first $S^2$ extends across the surgery
trace, since it evaluates to zero on the attaching sphere.
The original pair $(S^2\times S^4,u)$ spin-bounds
$S^2\times D^5$ over $K(\mathbb Z/3,2)$.
Thus the resulting model has zero spin-bordism class over this
Eilenberg--Mac Lane space. Zhubr's Theorem~3
\cite[p.795]{manZhubr75} identifies it, orientation-preservingly,
with every manifold in (D.T31), since both classifying maps induce
isomorphisms on $H_2$ and both have $b_3=0$.
\end{proof}
